\section{句法歧义：理解语言的核心挑战}

人类语言与编程语言的一个根本区别在于：\textbf{人类语言存在全局歧义（global ambiguity）}，而编程语言的歧义总能通过语法规则唯一消解。

\subsection{介词短语附着歧义（PP Attachment Ambiguity）}

\begin{figure}[H]
\centering
\includegraphics[width=0.95\textwidth]{slides-images/slide-006.jpg}
\caption{``Scientists count whales from space'' 的两种解读\protect\footnotemark}
\end{figure}
\footnotetext{Slides 第6页。}

以新闻标题 ``Scientists count whales from space'' 为例：

\begin{itemize}
    \item \textbf{解读 1}：``from space'' 修饰 ``count'' —— 科学家\textbf{从太空}数鲸鱼（正确解读）
    \item \textbf{解读 2}：``from space'' 修饰 ``whales'' —— 科学家数\textbf{来自太空的}鲸鱼
\end{itemize}

这种歧义源于英语中介词短语（PP）可以附着到不同的中心词上。由于 PP 在英语中极为常见，这类歧义也极为普遍。

\begin{warningbox}{PP 附着歧义的指数爆炸}
当句子中有多个连续的介词短语时，可能的解读数量按\textbf{Catalan 数}增长（指数级）。Manning 举了 Wall Street Journal 的例子：``The board approved its acquisition by Royal Trustco Ltd. of Toronto for \$27 a share at its monthly meeting'' 包含 4 个连续 PP，理论上有 \textbf{14 种}可能的依存结构（实际由上下文和世界知识消歧）。
\end{warningbox}

\begin{figure}[H]
\centering
\includegraphics[width=0.95\textwidth]{slides-images/slide-007.jpg}
\caption{PP 附着歧义的两种依存结构：``from space'' 分别附着到动词和名词\protect\footnotemark}
\end{figure}
\footnotetext{Slides 第7页。}

\subsection{并列结构作用域歧义（Coordination Scope Ambiguity）}

\begin{figure}[H]
\centering
\includegraphics[width=0.95\textwidth]{slides-images/slide-009.jpg}
\caption{并列结构作用域歧义示例\protect\footnotemark}
\end{figure}
\footnotetext{Slides 第9页。}

另一类常见歧义来自\textbf{并列结构}（coordination）的作用域。例如：

``Shuttle veteran and longtime NASA executive Fred Gregory appointed to board''

\begin{itemize}
    \item \textbf{解读 1}：Fred Gregory 是一个人，既是 shuttle veteran 又是 NASA executive
    \item \textbf{解读 2}：有两个人——一个 shuttle veteran 和 NASA executive Fred Gregory
\end{itemize}

类似的歧义还出现在 ``Doctor: No heart, cognitive issues'' 中——``No'' 的作用域是仅覆盖 ``heart'' 还是同时覆盖 ``heart'' 和 ``cognitive issues''。

\begin{figure}[H]
\centering
\includegraphics[width=0.95\textwidth]{slides-images/slide-010.jpg}
\caption{``Doctor: No heart, cognitive issues'' 的两种解读\protect\footnotemark}
\end{figure}
\footnotetext{Slides 第10页。}

\subsection{修饰语附着歧义}

\begin{figure}[H]
\centering
\includegraphics[width=0.95\textwidth]{slides-images/slide-012.jpg}
\caption{``Mutilated body washes up on Rio beach to be used for Olympics beach volleyball'' 的歧义\protect\footnotemark}
\end{figure}
\footnotetext{Slides 第12页。}

不定式短语 ``to be used for Olympics beach volleyball'' 可以修饰 ``beach''（海滩将用于沙滩排球）或 ``body''（尸体将用于沙滩排球），产生截然不同的含义。

\begin{knowledgebox}{不同语言的歧义类型不同}
Manning 特别指出，不同语言的句法歧义类型并不完全相同。例如，中文的介词短语修饰动词时通常出现在动词之前，宾语在动词之后，因此不会出现英语中 PP 修饰动词 vs. 修饰宾语的歧义。但中文有自己独特的歧义类型。
\end{knowledgebox}

\subsection{歧义的实际影响}

Catalan 数列决定了 PP 附着歧义的组合爆炸速度：

\begin{table}[H]
\centering
\begin{tabular}{cc}
\toprule
\textbf{PP 数量} & \textbf{可能的解读数} \\
\midrule
1 & 2 \\
2 & 5 \\
3 & 14 \\
4 & 42 \\
5 & 132 \\
\bottomrule
\end{tabular}
\caption{PP 数量与可能解读数的关系（Catalan 数）}
\end{table}

尽管歧义数量指数增长，人类读者通常能毫不费力地选择正确解读——这依赖于世界知识、上下文和语用推理。NLP 系统需要从数据中学会类似的消歧能力。

\subsection{句法结构在 NLP 中的应用}

\begin{figure}[H]
\centering
\includegraphics[width=0.95\textwidth]{slides-images/slide-013.jpg}
\caption{依存结构在信息抽取中的应用：从生物医学文本中提取蛋白质-蛋白质相互作用\protect\footnotemark}
\end{figure}
\footnotetext{Slides 第13页。}

依存结构在实际 NLP 系统中有广泛应用。例如在生物信息学领域，可以通过依存分析从文本中提取蛋白质-蛋白质相互作用事实。动词 ``interacts'' 的主语和介词宾语分别指向两个交互的蛋白质，并列结构则帮助识别多重交互关系。

\subsection{本章小结}

句法歧义是自然语言的根本特征，包括 PP 附着歧义、并列结构作用域歧义和修饰语附着歧义等。歧义数量随句子复杂度指数增长，但人类能利用世界知识和上下文轻松消歧。NLP 系统需要从标注数据中学习这种消歧能力。理解并正确解析句法结构对信息抽取、语义理解等下游任务至关重要。

%% ========== 依存语法 ==========

